L'isòtop $^{99\mathrm{m}}\mathrm{Tc}$, tecneci metaestable, s'utilitza com a radiotraçador en medicina, i té un període de semidesintegració de 6 hores, un temps suficient perquè s'acumuli en l'òrgan que es vol estudiar sense que perduri gaire temps en l'organisme. Aquest isòtop és un emissor $\gamma$ amb una energia d'uns 140~keV, i se'n pot generar una imatge. El tecneci és l'element 43 de la taula periòdica.

\begin{apartat}{1}
Escriviu el nombre màssic, el nombre atòmic i el nombre de neutrons que conté aquest isòtop. Escriviu la reacció de desintegració $\gamma$ que es produeix.
\end{apartat}

\begin{apartat}{1}
Si un pacient rep una dosi de 2~ng de l'isòtop, quina quantitat romandrà al seu cos passades 24 hores, si suposem que no n'ha eliminat gens per l'orina?
\end{apartat}
